% chktex-file 8
%   Allow dashes
% chktex-file 17
%   Don't require matching number of brackets and parentheses
% chktex-file 11
%   Allow ellipsis
% chktex-file 44
%   Allow regex
% chktex-file 9
%   Allow unmatched brackets and parentheses
% chktex-file 36
%   ALlow parentheses without spaces in front of them
% chktex-file 40
%   Allow factorials
% chktex-file 12
%   Don't require inter-word spaces
% chktex-file 3
%   Don't require enclosing parentheses with braces
% chktex-file 13
%   Don't require intersentence spaces
% chktex-file 26
%   Allow spaces in front of punctuation
\documentclass{article}
\usepackage{amsmath}
\usepackage{amssymb}
\usepackage[a4paper, margin=1in]{geometry}
\usepackage{graphicx}
\usepackage{hyperref}
\usepackage{listings}
\usepackage{cancel}
\lstset{
basicstyle=\small\ttfamily,
columns=flexible,
breaklines=true
}

\begin{document}

\section*{Construction Site Notes}
\begin{itemize}
    \item User flow:
    \begin{enumerate}
        \item Click on a zone
        \item Click button to start designing a building. Each building type can have its own button.
        \item Initiating design enters a new `blueprint mode', where you specify any desired details of the building - this creates a
            hypothetical building `blueprint' visible at the zone. Only one such blueprint can exist per zone at a time.
        \begin{itemize}
            \item Building details include building type, tolerance to environmental conditions, quality, etc., which determine materials
                required per construction step attempt, number of successful attempts required to complete, and expected value of step
                attempts and materials spent to complete
            \item New concept: \textbf{interface container}: Subclass of interface collection which is dynamically created
                and removed, contains a variety of interface, includes an intrinsic image with tooltip and boundaries, and
                can be moved around the screen. Throughout the game, this pattern should be re-used for temporary, complex
                configuration interfaces, like the blueprint mode, organism design, contract bids, complex events, etc.
        \end{itemize}
        \item Upon confirming details, a small amount of construction materials are immediately allocated, and a construction site is
            created. There should be another interface mode that describes any construction sites or completed buildings in the
            zone, which should be able to support multiple of these existing at once (scrolling through them). These specifications
            should reflect the player's current posterior belief about the building, but may not have actually been implemented 
            fully (update belief upon specific evidence arising).
        \item When a construction site and construction crew are selected, the construction crew can be attached to the site.
        \item During the end of the turn, checks should be rolled to make progress on the site, using up materials over time.
        \begin{itemize}
            \item It should be toggleable by a construction site whether these rolls will be shown to the player or automated. Whether they
                are shown should never affect the outcome.
            \item Minister theft can occur via either direct theft of materials, theft of wages of crew(s), and/or using a substandard
                amount of materials to create a slightly off-blueprint building and keeping the excess
        \end{itemize}
        \item Upon completion, the building is created where the site was, and the construction crew(s) are unattached.
    \end{enumerate}
    \item A pre-built module acts as the above, but selecting the pre-built module item opens up a selector to create a
        deployment site in some zone in the location, which opens a blueprint mode with pre-filled, immutable details, and with a
        very abbreviated construction process.
    \item Designing a contracted (or built on-site) module or spaceship, or an engineered microbe, or a variety of equipment should use a very similar system.
    \item Idea - potentially have about 6 progress steps required for each building, but each success can give some number of progress steps.
\end{itemize}

\end{document}

% Sample figure
%\begin{figure}[!ht]
%    \centering
%    \includegraphics[width=350px]{screenshot_1.png}
%    \caption{Sample caption}
%\end{figure}

% Sample hyperlink
% \item \href{URL}{Label}

% Sample listing
% \begin{lstlisting}[language=Python, caption=Sample Python Code]
% print("Hello World")
% \end{lstlisting}
